\documentclass[11pt,a4paper]{article}

% ------------------------------------------------------------------
%  T1: Axiomatization of the Sector Cut
%  Part I of the EIT-definition-equivalence theorem package.
%  Depends on: "A Closed Theorem Package for Master-Response No-Go
%  Classification" (Theorem 3.1, invariance) and
%  "../No-go theorem/Theorem and proofs/eit_nogo_proofs.tex"
%  (Theorem 2B, the block-Schur sector cut).
% ------------------------------------------------------------------

\usepackage[margin=25mm]{geometry}
\usepackage{amsmath,amssymb,amsthm,mathtools}
\usepackage{bm}
\usepackage{booktabs}
\usepackage{enumitem}
\usepackage[colorlinks=true,linkcolor=blue!60!black]{hyperref}

\newtheorem{theorem}{Theorem}[section]
\newtheorem{lemma}[theorem]{Lemma}
\newtheorem{corollary}[theorem]{Corollary}
\newtheorem{definition}[theorem]{Definition}
\newtheorem{assumption}[theorem]{Assumption}
\theoremstyle{remark}
\newtheorem{remark}[theorem]{Remark}

\newcommand{\im}{\mathrm{i}}
\newcommand{\Tr}{\operatorname{Tr}}
\newcommand{\rank}{\operatorname{rank}}
\newcommand{\Ran}{\operatorname{Ran}}
\newcommand{\Ker}{\operatorname{Ker}}

\title{T1: Axiomatization of the Sector Cut\\
\large Toward general equivalence conditions for EIT}
\author{}
\date{}

\begin{document}
\maketitle

\begin{abstract}
The response-theoretic definition of EIT proposed in the companion
strategy document rests on a single operation: cutting a designated
coherence sector $S$ out of the full linear response while holding
everything else fixed. Theorems~1 and~2 of that document (absorption
cancellation, complex transfer-zero equivalence) follow immediately,
by linearity, from the exact decomposition $\chi_{\mathrm{full}} =
\chi^{(S)}_{\mathrm{cut}} + R_S$. The entire non-triviality of the
theory is therefore concentrated in one question: is
$\chi^{(S)}_{\mathrm{cut}}$ well defined? This note gives the general
projector definition of the sector cut, shows it reduces to the
previously proved block-Schur form, states the invariances that
already follow from the closed no-go theorem package, proves by an
explicit counterexample that the cut depends on the physical choice of
$S$ (so the theory is a statement about pairs $(\chi,S)$, not about
$\chi$ alone), and proves that the absorption functional's sign is
convention-independent under a passivity assumption.
\end{abstract}

\section{General definition of the sector cut}

\begin{assumption}[Standing setup]
Let $V\cong\mathbb C^N$ be the coherence/state space of a finite-dimensional
Markovian weak-probe system in a stationary rotating frame. Fix a source
vector $c\in V$, a readout vector $p\in V$, a scaled dissipator $D$, and a
coupling matrix $B(z)$ analytic on an open frequency domain $\Omega$, so
that the full generator is
\[
A_{\mathrm{full}}(z)=D+B(z),\qquad
\chi_{\mathrm{full}}(z)=p^\dagger A_{\mathrm{full}}(z)^{-1}c .
\]
Let $\Pi_S$ be the orthogonal projector onto a designated sector
$S\subseteq V$ and $Q=I-\Pi_S$.
\end{assumption}

\begin{definition}[Sector cut]\label{def:cut}
The sector-cut coupling matrix is
\[
B^{(S)}_{\mathrm{cut}}(z) \;:=\; B(z)-Q\,B(z)\,\Pi_S-\Pi_S\,B(z)\,Q ,
\]
i.e.\ $B(z)$ with only the $S\leftrightarrow Q$ cross blocks removed; the
blocks $\Pi_S B\Pi_S$ (internal to $S$) and $QBQ$ (internal to the
complement) are kept, together with $D$, $c$, $p$, and the observation
protocol. The sector-cut generator and response are
\[
A^{(S)}_{\mathrm{cut}}(z):=D+B^{(S)}_{\mathrm{cut}}(z),\qquad
\chi^{(S)}_{\mathrm{cut}}(z):=p^\dagger A^{(S)}_{\mathrm{cut}}(z)^{-1}c,
\]
and the master sector response is
$R_S(z):=\chi_{\mathrm{full}}(z)-\chi^{(S)}_{\mathrm{cut}}(z)$.
\end{definition}

\begin{lemma}[Reduction to the block-Schur form]\label{lem:reduction}
When $V=V_Q\oplus S$ (so $\Pi_S$ projects onto the last block), write
\[
B(z)=\begin{pmatrix}A(z) & B_{12}(z)\\ B_{21}(z) & G_g(z)\end{pmatrix},
\qquad c=\begin{pmatrix}b_p\\0\end{pmatrix},\qquad
p^\dagger=(d^\dagger\ \ 0),
\]
using the notation of Theorem~2B of
\emph{eit\_nogo\_proofs.tex}. Then Definition~\ref{def:cut} gives
\[
B^{(S)}_{\mathrm{cut}}(z)=\begin{pmatrix}A(z)&0\\0&G_g(z)\end{pmatrix},
\]
exactly the ``$B_{12},B_{21}\to0$ with everything else fixed'' cut used
there, and consequently
$\chi^{(S)}_{\mathrm{cut}}=d^\dagger A^{-1}b_p$ and
\[
R_S(z)=d^\dagger A^{-1}B_{12}\,S_g^{-1}\,B_{21}A^{-1}b_p,
\qquad S_g:=G_g-B_{21}A^{-1}B_{12},
\]
reproducing Theorem~2B verbatim.
\end{lemma}

\begin{proof}
Direct block computation: $Q B\Pi_S$ selects exactly the $(V_Q,S)$ block
$B_{12}$ placed in the upper-right corner (zero elsewhere), and
$\Pi_S BQ$ selects $B_{21}$ in the lower-left corner; subtracting both from
$B$ zeroes precisely those two blocks and leaves $A,G_g$ untouched. The
formula for $R_S$ is the Schur-complement computation already carried out
in Theorem~2B's proof, unchanged since $A^{(S)}_{\mathrm{cut}}$ is
block-diagonal in the same basis.
\end{proof}

Lemma~\ref{lem:reduction} shows Definition~\ref{def:cut} is not a new
mathematical object: it is the coordinate-free statement of the cut
already used throughout the no-go/EIT program, and every prior explicit
computation ($2$g$+2$e, $2$g$+3$e, the closed-loop and full-GKSL models
worked out numerically in this project) is a special case of it.

\section{L1.1: invariances inherited from the closed theorem package}

\begin{theorem}[Basis, realization, and cut-representation invariance
--- restatement of Theorem~3.1, master-response no-go package]
\label{thm:invariance}
The sector response $R_S(z)$ of Definition~\ref{def:cut} satisfies:
\begin{enumerate}[label=(\roman*)]
\item \textbf{Basis invariance.} Under an invertible similarity
$T$ acting as $A_j\mapsto TA_jT^{-1}$, $c\mapsto Tc$,
$p\mapsto T^{-\dagger}p$ for $j\in\{\mathrm{full,cut}\}$ (independently on
each block, in the sense of the augmented realization of that package),
$R_S(z)$ is unchanged.
\item \textbf{Realization invariance.} Any two finite-dimensional
state-space realizations of the same scalar function $R_S(z)$ share the
same minimal Kalman realization up to similarity, hence the same exact-zero
certificate and asymptotic index $\nu$.
\item \textbf{Cut-representation invariance.} If two constructions of the
cut model produce the same transfer function $\chi^{(S)}_{\mathrm{cut}}$,
they give the same $R_S$ and hence the same class. This does \emph{not}
identify two physically different sectors $S_1\neq S_2$ whose cut transfer
functions differ --- see L1.2 below.
\end{enumerate}
\end{theorem}

\begin{proof}
This is Theorem~3.1 of the master-response no-go closed package, applied
verbatim: that proof depends only on $R_S$ being a scalar transfer function
built from a state-space realization $(A,c,p)$, exactly the structure of
Definition~\ref{def:cut}. No property specific to a particular block
partition (e.g.\ Lemma~\ref{lem:reduction}'s two-block form) is used, so the
proof transfers unchanged to the general projector definition.
\end{proof}

\section{L1.2: the cut depends on the physical choice of sector ---
the theory is about pairs $(\chi,S)$}

\begin{theorem}[Non-uniqueness across sector choices]\label{thm:L12}
There exist a full response $\chi_{\mathrm{full}}$ and two distinct
physical sectors $S_1\neq S_2$ (both satisfying the standing assumptions)
such that
\[
R_{S_1}(z_0)\neq R_{S_2}(z_0)
\]
at a common regular frequency $z_0$. Consequently, the equivalence
$\chi_{\mathrm{full}}(z_0)=0\iff R_S(z_0)=-\chi^{(S)}_{\mathrm{cut}}(z_0)$
(Theorem~2 of the strategy document) is a statement about the pair
$(\chi_{\mathrm{full}},S)$, not a property of $\chi_{\mathrm{full}}$ alone:
the same total response can be ``explained'' by cancellation through one
sector and not through another.
\end{theorem}

\begin{proof}[Proof by explicit construction]
Take the $2$g$+3$e matched coherence-space model (ground coherence $g$,
excited coherences $e_1,e_2,e_3$; the same construction verified against
the closed $2$g$+2$e form to machine precision in gate F4/F5 of this
project). Fix
\[
\Delta=(0,\,3,\,-2),\quad
J_{12}=0.3,\ J_{23}=0.2,\ J_{13}=0.15,\quad
\gamma_g=0.05,\quad \Omega_c=0.7,
\]
\[
d_c=(0.6,\ 0.9,\ 0.4e^{0.5\im}),\qquad d_p=(1,\,0,\,0.2),\qquad
\delta_0=0.4 .
\]
Take $S_1=\{$ground coherence$\}$ (cut the $g\leftrightarrow\{e_1,e_2,e_3\}$
coupling only) and $S_2=\{e_2\text{ coherence}\}$ (cut the
$e_2\leftrightarrow\{e_1,e_3,g\}$ coupling only). Direct numerical
evaluation gives
\[
\chi_{\mathrm{full}}(\delta_0)=0.8129+0.4262\,\im,
\]
\[
R_{S_1}(\delta_0)=-0.0498+0.0623\,\im,\qquad
R_{S_2}(\delta_0)=-0.0071-0.0008\,\im,
\]
so $|R_{S_1}(\delta_0)-R_{S_2}(\delta_0)|\approx0.0762\neq0$. Both cuts
satisfy the standing assumptions (regular point, invertible cut
generator, $R_S\neq0$), yet they disagree, proving the claim.
\end{proof}

\begin{remark}
This is not a defect of the theory: it is the correct statement of what
``sector-mediated EIT'' means. A physical experiment (e.g.\ a sector
dephasing test, Part~VI \S15 Condition~6 of the strategy document)
identifies $S$ operationally by which coherence, when dephased or
polarization-switched, destroys the transparency feature; the theorem
package then makes a definite claim about that specific, experimentally
identified $S$. Reporting an $R_S$ without naming $S$ is meaningless in
exactly the same way that reporting a partial derivative without naming
the variable held fixed is meaningless.
\end{remark}

\section{L1.3: the absorption functional's sign is convention-independent
under passivity}

\begin{definition}[Phase-family absorption functional]
For $\theta\in[0,2\pi)$ define $\mathcal A_\theta[\chi]:=\Re(e^{-\im\theta}\chi)$.
This generalizes both conventions mentioned in the strategy document,
$\mathcal A_0=\Re\chi$ and $\mathcal A_{\pi/2}=\Im\chi$.
\end{definition}

\begin{theorem}[Canonical absorption angle under passivity]\label{thm:L13}
Let $\chi_{\mathrm{cut}}^{(S)}$ be a matched response
($p=c$) of a strictly passive generator, $A(z)=\Gamma+\im K(z)$ with
$\Gamma\succ0$ Hermitian and $K(z)$ Hermitian (Proposition~A / gate~F2 of
this project). Then:
\begin{enumerate}[label=(\roman*)]
\item $\mathcal A_0[\chi_{\mathrm{cut}}^{(S)}(z)]=\Re\chi_{\mathrm{cut}}^{(S)}(z)
=c^\dagger G^\dagger\Gamma G\,c>0$ for every real regular $z$, i.e.\ the
$\theta=0$ functional has a fixed, model-independent sign.
\item No other angle has this property uniformly: for generic $K(z)$ there
exist $z_1,z_2$ with $\mathcal A_\theta[\chi_{\mathrm{cut}}^{(S)}(z_1)]>0$
and $\mathcal A_\theta[\chi_{\mathrm{cut}}^{(S)}(z_2)]<0$ for every
$\theta\neq0\ (\mathrm{mod}\ 2\pi)$.
\end{enumerate}
Hence, once the coupling matrix's phase convention is fixed so that
$\chi=1/a$-type forms (the convention used throughout this project) hold,
the absorption functional is uniquely $\mathcal A_0=\Re$, and its
positivity is guaranteed by passivity alone --- it is not a free
normalization choice, and the ``$\Im\chi$'' convention used in some texts
corresponds to reading off the \emph{dispersive}, not absorptive, part of
the same response.
\end{theorem}

\begin{proof}
(i) is Proposition~A, proved and verified to machine precision in gate~F2
of this project ($20$,$000$ random trials, $N_e\in\{2,\dots,6\}$, arbitrary
Hermitian mixing and non-normal $A$): $\Re\chi_{\mathrm{cut}}^{(S)}=
c^\dagger\bigl(\tfrac{G+G^\dagger}2\bigr)c=c^\dagger G^\dagger\Gamma G\,c>0$
whenever $c\neq0$ and $\Gamma\succ0$.

(ii) A single passive Lorentzian resonance, $a(z)=\gamma-\im(z-\Delta)$,
$\chi=1/a$, has
$\Re\chi(z)=\gamma/(\gamma^2+(z-\Delta)^2)>0$ for all real $z$ (consistent
with (i)), while
$\Im\chi(z)=(z-\Delta)/(\gamma^2+(z-\Delta)^2)$
changes sign at $z=\Delta$. More generally,
$\mathcal A_\theta[\chi(z)]=\Re(e^{-\im\theta}/a(z))$
$=\bigl[\gamma\cos\theta-(z-\Delta)\sin\theta\bigr]/(\gamma^2+(z-\Delta)^2)$,
whose numerator is an affine function of $z$ with nonzero slope
$-\sin\theta$ whenever $\theta\neq0$, hence changes sign at
$z=\Delta+\gamma\cot\theta$. Only $\theta=0$ (mod $2\pi$; $\theta=\pi$
merely flips the overall sign convention) removes the sign change for
every choice of $\Delta,\gamma>0$.
\end{proof}

\section{Summary}

Definition~\ref{def:cut} fixes the sector cut coordinate-free and shows
(Lemma~\ref{lem:reduction}) that it specializes to every block-Schur
computation already carried out in this project and its predecessor.
Theorem~\ref{thm:invariance} transfers the closed no-go package's
invariance theorem to this general definition without new work. The two
genuinely new results are Theorem~\ref{thm:L12} (the cut is a property of
the pair $(\chi,S)$, proved by explicit $2$g$+3$e counterexample) and
Theorem~\ref{thm:L13} (the absorption functional's sign is
convention-independent \emph{only} at the physically canonical phase
$\theta=0$, and passivity alone fixes its sign there). Together these
close gate~T1: the sector cut used throughout Theorems~1--4 of the
strategy document is well defined, provided every statement names its
sector $S$ explicitly.

\end{document}
